\corrigesChapitre{8}{Décrire et mesurer un mouvement}
\par\noindent\begin{minipage}{\linewidth}
\begin{corrige}{1}{Le référentiel}
\begin{questions}
\item Oui, s’il reste posé à la même place.
\item Oui lorsque le train roule ; la réponse dépend du référentiel.
\end{questions}
\pcFigureEx{figures/cor-referentiels.pdf}{Le mouvement d’un même objet peut être décrit différemment selon le référentiel.}
\end{corrige}
\end{minipage}
\par\noindent\begin{minipage}{\linewidth}
\begin{corrige}{2}{La trajectoire}
\begin{questions}
\item Circulaire.
\item Non : sa trajectoire est rectiligne dans le référentiel considéré.
\end{questions}
\end{corrige}
\end{minipage}
\par\noindent\begin{minipage}{\linewidth}
\begin{corrige}{3}{Temps égaux}
\begin{questions}
\item Elle augmente : le mouvement est accéléré au sens de sa valeur de vitesse.
\item Les intervalles de temps doivent être réellement égaux et les longueurs comparables sur l’image.
\end{questions}
\pcFigureEx{figures/cor-chronophoto.pdf}{Comparer les distances entre positions successives, mesurées pour des durées égales.}
\end{corrige}
\end{minipage}
\par\noindent\begin{minipage}{\linewidth}
\begin{corrige}{4}{Une vitesse directe}
\begin{questions}
\item $v=40/8=5\unite{m/s}$.
\item Mètre par seconde, car les données sont en mètres et secondes.
\end{questions}
\par\smallskip\noindent\textbf{Pour comprendre.}\enspace Choisir $v=d/\Delta t$ avec d en m et la durée en s : $v=40/8=5\unite{m/s}$. Il s’agit de la moyenne sur le trajet, sauf information supplémentaire.\par
\end{corrige}
\end{minipage}
\par\noindent\begin{minipage}{\linewidth}
\begin{corrige}{5}{Convertir}
\begin{questions}
\item $10\times3{,}6=36\unite{km/h}$.
\item $72/3{,}6=20\unite{m/s}$.
\end{questions}
\par\smallskip\noindent\textbf{Pour comprendre.}\enspace $1\unite{m/s}=3{,}6\unite{km/h}$. Multiplier par 3,6 vers km/h et diviser vers m/s : 36 km/h et 20 m/s.\par
\end{corrige}
\end{minipage}
\par\noindent\begin{minipage}{\linewidth}
\begin{corrige}{6}{Calculer une distance}
\begin{questions}
\item $d=0{,}40\times15=6{,}0\unite m$.
\item $12\unite m$ si la vitesse reste la même.
\end{questions}
\end{corrige}
\end{minipage}
\par\noindent\begin{minipage}{\linewidth}
\begin{corrige}{7}{Calculer une durée}
\begin{questions}
\item $\Delta t=1{,}8/6=0{,}30\unite h$.
\item $0{,}30\times60=18\unite{min}$.
\end{questions}
\par\smallskip\noindent\textbf{Pour comprendre.}\enspace Pour une vitesse moyenne fournie, $\Delta t=d/v=1{,}8/6=0{,}30\unite h$. Puis $0{,}30\times60=18\unite{min}$ ; 0,30 h ne signifie pas 30 min.\par
\end{corrige}
\end{minipage}
\par\noindent\begin{minipage}{\linewidth}
\begin{corrige}{8}{Lire des positions}
\begin{questions}
\item Les écarts sont toujours de $3\unite m$ en $2\unite s$ : la valeur de la vitesse est constante.
\item $v=3/2=1{,}5\unite{m/s}$.
\end{questions}
\pcFigureEx{figures/cor-position.pdf}{Des distances de 3 m sont parcourues en 2 s ; la pente traduit une vitesse constante de 1,5 m/s.}
\end{corrige}
\end{minipage}
\par\noindent\begin{minipage}{\linewidth}
\begin{corrige}{9}{Uniforme sur un cercle}
\begin{questions}
\item Non : la direction est tangente au cercle et varie avec la position.
\item C’est un mouvement circulaire uniforme ; le vecteur vitesse n’est pas constant.
\end{questions}
\pcFigureEx{figures/cor-manege.pdf}{Les flèches de vitesse sont tangentes au cercle : leur direction change même si leur valeur reste la même.}
\end{corrige}
\end{minipage}
\par\noindent\begin{minipage}{\linewidth}
\begin{corrige}{10}{Une vidéo étalonnée}
\begin{questions}
\item $80/200=0{,}40\unite m$, si règle et mouvement sont dans le même plan.
\item $0{,}40/0{,}20=2{,}0\unite{m/s}$.
\end{questions}
\par\smallskip\noindent\textbf{Pour comprendre.}\enspace La calibration donne $1/200=0{,}005\unite m$ par pixel. Donc $d=80\times0{,}005=0{,}40\unite m$, puis $v=0{,}40/0{,}20=2{,}0\unite{m/s}$.\par
\pcFigureEx{figures/rev-cor-pixels.pdf}{Convertir le déplacement de l’image en distance réelle avant de diviser par la durée.}
\end{corrige}
\end{minipage}
\par\noindent\begin{minipage}{\linewidth}
\begin{corrige}{11}{Deux parties d’un trajet}
\begin{questions}
\item $d_{\rm total}=120\unite{km}$ et $\Delta t=3\unite h$ ; $v=40\unite{km/h}$.
\item $45$ est la moyenne arithmétique de 60 et 30, mais les durées ne sont pas égales.
\end{questions}
\par\smallskip\noindent\textbf{Pour comprendre.}\enspace Additionner d’abord les distances et les durées : $d=120\unite{km}$, $\Delta t=3\unite h$, donc $v_{\rm moy}=40\unite{km/h}$. La moyenne arithmétique des deux vitesses ne convient pas lorsque les durées diffèrent.\par
\end{corrige}
\end{minipage}
\par\noindent\begin{minipage}{\linewidth}
\begin{corrige}{12}{Un arrêt compte-t-il ?}
\begin{questions}
\item $30\unite{min}=0{,}5\unite h$ ; $v=12/0{,}5=24\unite{km/h}$.
\item Non : pendant l’arrêt, la vitesse par rapport au sol est nulle.
\end{questions}
\par\smallskip\noindent\textbf{Pour comprendre.}\enspace Le trajet complet dure $20+10=30\unite{min}=0{,}50\unite h$. La distance reste 12 km : $v_{\rm moy}=12/0{,}50=24\unite{km/h}$.\par
\pcFigureEx{figures/rev-cor-bus.pdf}{Pour la vitesse moyenne sur le trajet complet, l’arrêt fait partie de la durée totale.}
\end{corrige}
\end{minipage}
\par\noindent\begin{minipage}{\linewidth}
\begin{corrige}{13}{Une conclusion trop forte}
\begin{questions}
\item Non : elle peut avoir accéléré puis ralenti.
\item Il donne une vitesse moyenne de 20 m/s sur ce tronçon, pas toute son évolution temporelle.
\end{questions}
\end{corrige}
\end{minipage}
\par\noindent\begin{minipage}{\linewidth}
\begin{corrige}{14}{Un graphique ambigu}
\begin{questions}
\item Non : une position décroissante peut correspondre à un mouvement uniforme dans le sens opposé à l’axe choisi.
\item Il faut distinguer signe de l’évolution de position et valeur positive de la vitesse.
\end{questions}
\par\smallskip\noindent\textbf{Pour comprendre.}\enspace La pente est négative parce que le mobile se déplace vers des positions plus petites. La droite a une pente de valeur constante : dans cet exemple rectiligne, cela ne traduit pas un ralentissement.\par
\pcFigureEx{figures/cor-position-retour.pdf}{Une position décroissante indique un sens de déplacement ; la valeur de la vitesse peut rester constante.}
\end{corrige}
\end{minipage}
\par\noindent\begin{minipage}{\linewidth}
\begin{corrige}{15}{Comparer des mesures}
\begin{questions}
\item $\bar t=(4{,}8+5{,}0+5{,}2)/3=5{,}0\unite s$.
\item $v\approx10/5{,}0=2{,}0\unite{m/s}$. La dispersion et le déclenchement manuel ne justifient pas de nombreuses décimales.
\end{questions}
\end{corrige}
\end{minipage}
\par\noindent\begin{minipage}{\linewidth}
\phantomsection\label{cor-defi-8}\pcSolutionDefi{8}{Choisir le meilleur trajet}\begin{questions}
\item A : $1{,}2/4=0{,}30\unite h=18\unite{min}$. B : $1{,}5/6=0{,}25\unite h=15\unite{min}$, puis 3 minutes : 18 minutes.
\item Les durées sont égales dans ce modèle.
\item Les vitesses moyennes, la durée d’arrêt ou les distances peuvent varier. La conclusion est conditionnelle aux données.
\end{questions}
\pcFigureEx{figures/rev-cor-trajets.pdf}{Inclure les trois minutes d’arrêt dans la durée du trajet B ; les deux arrivées sont simultanées dans ce modèle.}
\end{minipage}
\par\noindent\begin{minipage}{\linewidth}
\phantomsection\label{cor-activite-8}\pcSousTitre{Pistes pour l’investigation}
Des distances successives sensiblement égales pour des durées égales sont compatibles avec une vitesse de valeur constante. Le graphe de position est alors proche d’une droite sur le déplacement rectiligne étudié. Contrôler l’échelle, la perspective et le même point suivi. Des écarts plus petits que la lecture ne justifient pas une conclusion certaine d’accélération.\par\end{minipage}
